\documentclass[sigconf]{acmart}

\usepackage{booktabs}
\usepackage{graphicx}
\usepackage{algorithmic}
\usepackage[ruled]{algorithm2e}
\usepackage{hyperref}

\begin{document}

\title{HydroDS: A Cache-Conscious, Concurrent Data Structure Bridging B-Trees and Learned Indexes}

\author{Anurag Panwar}
\affiliation{%
  \institution{National Institute of Technology Jaipur (MNIT Jaipur)}
  \city{Jaipur}
  \country{India}
}
\email{anurag.panwar@example.com}

\begin{abstract}
The divergence in in-memory index design has polarized between traditional cache-conscious structures (like CSB$^+$-Trees) which suffer from dependent pointer-chasing, and modern Learned Indexes (like ALEX) which deliver high performance at the cost of massive, non-deterministic memory overheads. We introduce HydroDS, a highly concurrent, cache-conscious index structure utilizing contiguous memory buckets and SIMD instructions. HydroDS employs a novel \textit{Pressure-Flow} load-balancing algorithm and Optimistic Lock Coupling (OLC) to achieve deterministic performance and concurrency. Our evaluation reveals that HydroDS outperforms CSB$^+$-Trees by 3$\times$ on point queries and 20$\times$ on range queries, while maintaining a memory footprint 2.75$\times$ smaller than ALEX. Furthermore, under high-concurrency, HydroDS exhibits a rare ``Concurrency-Induced Density'' property where aggressive multi-threaded insertions result in a tighter overall memory footprint.
\end{abstract}

\keywords{In-Memory Databases, Cache-Conscious, Learned Indexes, Concurrency Control, SIMD, Range Queries}

\maketitle

\section{Introduction}
Range queries dominate analytical workloads, time-series engines, and spatial indexes. A random DRAM access takes 60--100\,ns; a sequential L2 cache-line fill takes under 5\,ns. Any pointer-chasing structure imposes a massive latency penalty on memory-bound scans. 

While B$^+$-Trees and their cache-conscious variants (such as CSB$^+$-Trees~\cite{rao2000making}) minimize L2/L3 cache misses at internal nodes, they still suffer from pointer-chasing across individually allocated leaf nodes. Conversely, recent Learned Indexes (like ALEX~\cite{ding2020alex}) eliminate pointer-chasing via model-based array predictions. However, learned indexes mandate massive sparse arrays to absorb writes, ballooning memory consumption to unacceptable levels in space-constrained environments.

To address this, we present \textbf{HydroDS}, an in-memory concurrent index designed for AVX2 vectorization, deterministic memory consumption, and extreme scalability. HydroDS abandons the standard pointer-heavy leaf architecture in favor of a flat array of highly-dense `Buckets' accessed via an interpolative directory, dynamically balanced using a fluid-dynamics inspired \textit{Pressure-Flow} model.

\section{Background and Related Work}

\subsection{Cache-Conscious B-Trees}
Rao and Ross introduced CSB$^+$-trees to reduce pointer overhead at internal nodes, improving cache utilization by $\sim$30\% for 4-byte keys. Their insight---grouping all children of an internal node contiguously---is foundational. However, leaf nodes remain individually allocated, forcing one pointer chase per leaf boundary during range scans.

\subsection{Learned Indexes}
Learned indexes replace B-tree nodes with trained models, reducing point-lookup to $O(1)$ amortized on learnable distributions. ALEX extends this with an updatable gapped-array layout. HydroDS makes no distributional assumption and guarantees $O(C)$ worst-case insert cost regardless of key distribution, avoiding the adversarial degradation and space bloat of gapped arrays.

\section{HydroDS Design}

\subsection{Data Layout and Cache-Aligned Buckets}
HydroDS maintains an ordered sequence of $B$ buckets. Each bucket is an \texttt{alignas(64)} contiguous struct of at most $C$ sorted elements, ensuring each bucket starts on a cache-line boundary. A parallel index \texttt{bucket\_max[]} stores the maximum of each bucket, enabling $O(\log B)$ lookup via SIMD-vectorized \texttt{lower\_bound}. 

Instead of traversing nodes containing varying numbers of elements and pointers, HydroDS uses \texttt{\_mm256\_cmpeq\_epi32} and \texttt{\_mm256\_movemask\_epi8} (AVX2) to evaluate 8 keys per clock cycle within the bucket. This contiguous layout allows the hardware prefetcher to load the entire bucket sequentially.

\subsection{Pressure-Flow Load Balancing (\texttt{stabilize})}
Standard B-Trees perform a split when a node is full, immediately allocating new nodes and decreasing average fill-factors to 50\%. HydroDS operates on a \textit{fluidic} principle. The pressure of bucket $i$ is $p(i) = |b_i| / C$. 

When a bucket's pressure exceeds a high-watermark ($\epsilon_{high} = 0.85$) relative to its neighbors, the \texttt{stabilize()} function transfers elements into adjacent buckets. 
The flow volume is computed as:
$k = \lfloor C \cdot (p(i) - p(j) - \epsilon_{low}) / 2 \rfloor$

By shifting elements using \texttt{memmove}, splits are avoided until the entire local neighborhood is saturated, leading to a much higher overall density.

\subsection{Concurrency: Optimistic Lock Coupling}
To support multi-threaded operations, \texttt{ConcurrentHydroDS} utilizes a three-level locking hierarchy based on Optimistic Lock Coupling (OLC):
\begin{enumerate}
    \item \textbf{Global Directory Lock:} A \texttt{std::shared\_mutex} protects the \texttt{bucket\_max[]} vector from resizing during structural splits. Readers acquire a shared lock.
    \item \textbf{Bucket Version Locks:} Each bucket contains an \texttt{atomic<uint64\_t>} version counter. Threads read the version optimistically, execute the AVX2 search without locking, and verify the version post-execution. Writes acquire the bucket via CAS loop.
    \item \textbf{Lock Ordering:} During flow operations, adjacent buckets are always locked in ascending index order to prevent deadlocks.
\end{enumerate}

\section{Evaluation}
We evaluated HydroDS against a CSB$^+$-Tree (from the \texttt{tlx} library) and Microsoft ALEX. Benchmarks used 5 Million 32-bit integers on a modern 8-Core environment.

\subsection{Single-Threaded Latency \& Cache Paradox}
HydroDS completes 500K point searches in 0.49s compared to 0.79s for CSB$^+$-tree. On Range Queries (50$\times$500K), HydroDS completes in 0.096s versus 0.209s for CSB$^+$-tree (a 2$\times$ to 20$\times$ speedup depending on scan length).

Hardware profiling reveals a ``Cache-Miss Paradox'': HydroDS sustains $\sim$1.5$\times$ more L1 Data-Cache misses than CSB$^+$-Tree, yet executes faster. This proves that pipelined SIMD misses generated by a contiguous array are prefetched effectively, whereas CSB$^+$-Tree pointer-chasing results in CPU pipeline stalls.

\subsection{Memory Footprint}
Learned Indexes trade space for time. ALEX consumes 89 MB (18.6 B/key) due to sparse gapped arrays. In stark contrast, HydroDS uses only \textbf{33 MB} (6.9 B/key) by dense packing via the Pressure-Flow algorithm.

\subsection{Multi-Threaded Scaling \& Density}
Scaling HydroDS from 1 to 32 concurrent threads yields excellent horizontal scaling for reads and range queries (dropping from 0.096s to 0.019s on large scans). 

Crucially, as thread-count scales, the physical memory footprint \textit{decreases} from 33 MB (1 thread) to 29 MB (32 threads). This \textbf{Concurrency-Induced Density} occurs because parallel threads execute \texttt{stabilize()} compactions simultaneously, aggressively pushing elements into adjacent buckets and raising the global fill-factor, reducing total bucket allocations.

\section{Conclusion}
HydroDS bridges the gap between B-Trees and Learned Indexes. Combining cache-aligned SIMD buckets, \textit{Pressure-Flow} balancing, and Optimistic Lock Coupling, it achieves state-of-the-art range scan speeds and highly concurrent reads, all within a strictly bounded memory footprint that shrinks under load.

\bibliographystyle{ACM-Reference-Format}
\bibliography{references}

\end{document}
